\Needspace{16\baselineskip}
\fancyhead[L]{\small VESTA.}
\begin{center}
VESTA.

\par\vspace{4pt}\hrule\vspace{4pt}

\fontsize{9}{11}\selectfont Göttingische gelehrte Anzeigen. Issue 55. Pp. 537..540. 8 April 1809.

\par\vspace{4pt}\hrule\vspace{4pt}
\end{center}

In earlier issues of our Gel. Anz. we reported the first observations of Vesta made at our observatory from June through September of last year, as well as those made in Berlin by Professor \textsc{Bode} (1808, p. 1793) [Nov. 10, p. 308 of this volume]. Professor \textsc{Gauss} recently received from Mr. \SourceLeaf{315}\textsc{Bouvard} the observations of this planet made at the Imperial Observatory in Paris. We believe that making them known early will be of service to astronomers.

\begin{center}
Observations of Vesta at Paris:
\end{center}


\begin{center}\fontsize{10}{12}\selectfont\setlength{\tabcolsep}{5pt}\renewcommand{\arraystretch}{1.15}\begin{tabular}{lrrr}
1808 & \shortstack{Mean time\\at Paris} & \shortstack{Apparent right\\ascension} & \shortstack{Apparent south\\declination}\\
\hline
Aug. 26 & $13^{\text{u}}\;16^{\text{m}}\;23^{\text{s}}5$ & $354^\circ\;11'\;28''80$ & $14^\circ\;1'\;20''3$\\
28 & $13\;6\;56.1$ & $353\;47\;30.30$ & $14\;17\;43.4$\\
Sept. 16 & $11\;35\;23.2$ & $349\;34\;9.45$ & $16\;34\;8.7$\\
19 & $11\;20\;57.3$ & $348\;54\;30.00$ & $16\;50\;16.0$\\
22 & $11\;6\;37.0$ & $348\;16\;15.00$ & $17\;4\;24.0$\\
24 & $10\;57\;7.7$ & $347\;51\;49.80$ & $17\;12\;48.5$\\
25 & $10\;52\;24.5$ & $347\;39\;55.35$ & $17\;16\;36.0$\\
26 & $10\;47\;42.1$ & $347\;28\;18.75$ & $17\;20\;10.0$\\
Oct.  4 & $10\;10\;48.6$ & $346\;6\;32.55$ & $17\;39\;50.5$\\
5 & $10\;6\;18.0$ & $345\;57\;50.25$ & $17\;41\;6.7$\\
6 & $10\;1\;48.7$ & $345\;49\;27.75$ & $17\;42\;20.0$\\
13 & $9\;31\;7.4$ & $345\;1\;51.75$ & $17\;43\;36.3$\\
\end{tabular}\end{center}


{\rightskip=0pt plus 1em \par\indent From these and the other available observations, Professor \textsc{Gauss} determined Vesta’s opposition, the \emph{first} observed so far.\par}

\Needspace{7\baselineskip}
\begin{center}Opposition of Vesta:

1808 September 8. $7^{\text{u}}\;52^{\text{m}}\;57^{\text{s}}$ Mean time at Göttingen.
\end{center}
\begin{center}\fontsize{10}{12}\selectfont\setlength{\tabcolsep}{5pt}\renewcommand{\arraystretch}{1.15}\begin{tabular}{lr}

True longitude .\;.\;.\;.\;.\;.\;.\;.\;.\; & $345^\circ\;53'\;38''5$\\
True geocentric latitude .\;.\;.\;.\; & $11\;0\;25.8$ south\\
\end{tabular}\end{center}


\par\indent At our observatory, meridian observations of Vesta continued until November and were then discontinued. As the planet grew fainter, the small aperture of the telescope on our mural quadrant no longer allowed a very accurate determination. Less accurate observations from that time would, in any case, have been of little use in improving the elements. For this reason, observations with the ring micrometer resumed only toward the end of this period of visibility. The unsettled weather allowed only two observations, but they turned out particularly well and are all the more worth reporting here because they were probably the last ones made during this period.


\begin{center}\fontsize{10}{12}\selectfont\setlength{\tabcolsep}{5pt}\renewcommand{\arraystretch}{1.15}\begin{tabular}{lrrr}
1809 & \shortstack{Mean time\\at Göttingen} & \shortstack{Apparent right\\ascension of Vesta} & \shortstack{Apparent south\\declination}\\
\hline
Feb. 8 & $7^{\text{u}}\;8^{\text{m}}\;25^{\text{s}}$ & $9^\circ\;26'\;51''3$ & $2^\circ\;9'\;28''8$\\
16 & $7\;16\;56$ & $12\;25\;20.0$ & $0\;41\;32.4$\\
\end{tabular}\end{center}


\par\indent The elements printed in issue 107 of last year’s Anzeigen [1808. July 4. p. 302 of this volume] differ from these latest observations by 4 minutes in longitude and only a few seconds in latitude. Those elements therefore required only a moderate adjustment to bring them into agreement with the observations of 1807, 1808, and 1809. Here are the new elements recently determined by Professor \textsc{Gauss}.



\Needspace{20\baselineskip}
\SourceLeaf{316}\begin{center}
IV. \emph{Elements of Vesta.}
\end{center}


\begin{center}\fontsize{10}{12}\selectfont\setlength{\tabcolsep}{5pt}\renewcommand{\arraystretch}{1.15}\begin{tabular}{lr}
\multicolumn{2}{c}{Epochs of mean longitude.}\\
\multicolumn{2}{c}{Meridian of Göttingen.}\\
1807 .\;.\;.\;.\;.\;.\;.\;.\;.\; & $168^\circ\;16'\;35''5$\\
1808 .\;.\;.\;.\;.\;.\;.\;.\;.\; & $267\;39\;28.6$\\
1809 .\;.\;.\;.\;.\;.\;.\;.\;.\; & $6\;46\;4.2$\\
1810 .\;.\;.\;.\;.\;.\;.\;.\;.\; & $105\;52\;39.7$\\
1811 .\;.\;.\;.\;.\;.\;.\;.\;\textsuperscript{.}\; & $204\;59\;15.3$\\
\end{tabular}\end{center}

\begin{center}\fontsize{9}{11}\selectfont\setlength{\tabcolsep}{5pt}\renewcommand{\arraystretch}{1.15}\begin{tabular}{lr}

Perihelion 1807, fixed with respect to the stars .\; & $249^\circ\;52'\;23''8$\\
Ascending node 1807, fixed with respect to the stars  & $103\;13\;11{,}2$\\
Inclination of the orbit .\;.\;.\;.\;.\;.\; & $7\;8\;18{,}8$\\
Daily mean tropical motion .\;.\; & $977''5221$\\
Tropical orbital period .\;.\; & 1325 days 19 hours\\
Eccentricity .\;.\;.\;.\;.\;.\;.\;.\;.\;.\; & $0{,}0887809$\\
Logarithm of the semimajor axis .\;.\; & $0{,}3732940$\\
\end{tabular}\end{center}


\par\indent According to these elements, we can expect Vesta’s second opposition on the afternoon of 1 January 1810 at longitude $100^\circ\;43'$ and geocentric latitude $0^\circ\;31'$ south.

\Needspace{18\baselineskip}
\par\indent Some of our readers may also be interested in the results on the varying relative brightness of this planet that Professor \textsc{Gauss} worked out in passing during these calculations. Taking as the unit the brightness the planet would show at distance 1 from Earth and Sun, its brightness on the day of discovery, 29 March 1807, when it looked like a star of magnitude 5..6, was $= 0,11815$. It was already decreasing then and had fallen to 0.02699 by 26 September 1807, when the Milan astronomers last observed it. On 22 June 1808, when the planet was first found again at our observatory, looking like a star of magnitude 8 or 7..8, its brightness was already back to $= 0.05492$. It grew brighter and brighter, and in September, when its brightness at opposition had risen to 0,09513, it looked like a star of magnitude 6 or 6..7. Its brightness then decreased again, and by 16 February it had fallen to 0,01629, though the planet was still easily as bright as a star of magnitude 9. At the next opposition it will rise only to 0,06435, so we can expect the planet to shine only like a star of magnitude 7.
